\chapter*{练习}
\addcontentsline{toc}{section}{练习}

\section*{练习 3.1}
\addcontentsline{toc}{subsection}{练习 3.1}

证明：对于所有 $ \boldsymbol{x} = [x_1, x_2]^\top \in \mathbb{R}^2 $ 和 $ \boldsymbol{y} = [y_1, y_2]^\top \in \mathbb{R}^2 $，由下式定义的 $ \langle\cdot, \cdot\rangle $：

\begin{equation*}
\langle\boldsymbol{x}, \boldsymbol{y}\rangle := x_1 y_1 - (x_1 y_2 + x_2 y_1) + 2(x_2 y_2)
\end{equation*}

是一个内积。

\medskip\hrule\medskip

\section*{练习 3.2}
\addcontentsline{toc}{subsection}{练习 3.2}

考虑 $ \mathbb{R}^2 $，其中对于 $ \mathbb{R}^2 $ 中的所有 $ \boldsymbol{x} $ 和 $ \boldsymbol{y} $，定义 $ \langle\cdot, \cdot\rangle $ 为：

\begin{equation*}
\langle\boldsymbol{x}, \boldsymbol{y}\rangle := \boldsymbol{x}^\top \underbrace{\begin{bmatrix} 2 & 0 \\ 1 & 2 \end{bmatrix}}_{=: \boldsymbol{A}} \boldsymbol{y}
\end{equation*}

$ \langle\cdot, \cdot\rangle $ 是一个内积吗？

\medskip\hrule\medskip

\section*{练习 3.3}
\addcontentsline{toc}{subsection}{练习 3.3}

计算

\begin{equation*}
\boldsymbol{x} = \begin{bmatrix} 1 \\ 2 \\ 3 \end{bmatrix}, \quad \boldsymbol{y} = \begin{bmatrix} -1 \\ -1 \\ 0 \end{bmatrix}
\end{equation*}

之间的距离，分别使用：

a. $ \langle\boldsymbol{x}, \boldsymbol{y}\rangle := \boldsymbol{x}^\top \boldsymbol{y} $ b. $ \langle\boldsymbol{x}, \boldsymbol{y}\rangle := \boldsymbol{x}^\top \boldsymbol{A} \boldsymbol{y} $，其中

\begin{equation*}
\boldsymbol{A} := \begin{bmatrix} 2 & 1 & 0 \\ 1 & 3 & -1 \\ 0 & -1 & 2 \end{bmatrix}
\end{equation*}

\medskip\hrule\medskip

\section*{练习 3.4}
\addcontentsline{toc}{subsection}{练习 3.4}

计算

\begin{equation*}
\boldsymbol{x} = \begin{bmatrix} 1 \\ 2 \end{bmatrix}, \quad \boldsymbol{y} = \begin{bmatrix} -1 \\ -1 \end{bmatrix}
\end{equation*}

之间的夹角，分别使用：

a. $ \langle\boldsymbol{x}, \boldsymbol{y}\rangle := \boldsymbol{x}^\top \boldsymbol{y} $ b. $ \langle\boldsymbol{x}, \boldsymbol{y}\rangle := \boldsymbol{x}^\top \boldsymbol{B} \boldsymbol{y} $，其中

\begin{equation*}
\boldsymbol{B} := \begin{bmatrix} 2 & 1 \\ 1 & 3 \end{bmatrix}
\end{equation*}

\medskip\hrule\medskip

\section*{练习 3.5}
\addcontentsline{toc}{subsection}{练习 3.5}

考虑带有点积的欧氏向量空间 $ \mathbb{R}^5 $。子空间 $ U \subseteq \mathbb{R}^5 $ 与向量 $ \boldsymbol{x} \in \mathbb{R}^5 $ 给定为：

\begin{equation*}
\begin{adjustbox}{max width=\linewidth}
U = \text{span}\left[ \begin{bmatrix} 0 \\ -1 \\ 2 \\ 0 \\ 2 \end{bmatrix}, \begin{bmatrix} 1 \\ -3 \\ 1 \\ -1 \\ 2 \end{bmatrix}, \begin{bmatrix} -3 \\ 4 \\ 1 \\ 2 \\ 1 \end{bmatrix}, \begin{bmatrix} -1 \\ -3 \\ 5 \\ 0 \\ 7 \end{bmatrix} \right], \quad \boldsymbol{x} = \begin{bmatrix} -1 \\ -9 \\ -1 \\ 4 \\ 1 \end{bmatrix}
\end{adjustbox}
\end{equation*}

a. 确定 $ \boldsymbol{x} $ 在 $ U $ 上的正交投影 $ \pi_U(\boldsymbol{x}) $。 b. 确定距离 $ d(\boldsymbol{x}, U) $。

\medskip\hrule\medskip

\section*{练习 3.6}
\addcontentsline{toc}{subsection}{练习 3.6}

考虑带有内积的 $ \mathbb{R}^3 $：

\begin{equation*}
\langle\boldsymbol{x}, \boldsymbol{y}\rangle := \boldsymbol{x}^\top \begin{bmatrix} 2 & 1 & 0 \\ 1 & 2 & -1 \\ 0 & -1 & 2 \end{bmatrix} \boldsymbol{y}
\end{equation*}

此外，定义 $ \boldsymbol{e}_1, \boldsymbol{e}_2, \boldsymbol{e}_3 $ 为 $ \mathbb{R}^3 $ 中的标准基。

a. 确定 $ \boldsymbol{e}_2 $ 在

\begin{equation*}
U = \text{span}[\boldsymbol{e}_1, \boldsymbol{e}_3]
\end{equation*}

上的正交投影 $ \pi_U(\boldsymbol{e}_2) $。 提示：正交性是通过该内积定义的。 b. 计算距离 $ d(\boldsymbol{e}_2, U) $。 c. 画出该场景示意图：标准基向量以及 $ \pi_U(\boldsymbol{e}_2) $。

\medskip\hrule\medskip

\section*{练习 3.7}
\addcontentsline{toc}{subsection}{练习 3.7}

设 $ V $ 为一个向量空间，$ \pi $ 为 $ V $ 上的一个自同态（endomorphism）。

a. 证明 $ \pi $ 是投影当且仅当 $ \text{id}_V - \pi $ 是投影，其中 $ \text{id}_V $ 是 $ V $ 上的恒等自同态。 b. 现在假设 $ \pi $ 是一个投影。计算 $ \text{Im}(\text{id}_V - \pi) $ 和 $ \text{ker}(\text{id}_V - \pi) $ 作为 $ \text{Im}(\pi) $ 和 $ \text{ker}(\pi) $ 的函数（即用 $ \text{Im}(\pi) $ 和 $ \text{ker}(\pi) $ 来表示它们）。

\medskip\hrule\medskip

\section*{练习 3.8}
\addcontentsline{toc}{subsection}{练习 3.8}

使用格拉姆-施密特方法（Gram-Schmidt method），将二维子空间 $ U \subseteq \mathbb{R}^3 $ 的一组基 $ B = (\boldsymbol{b}_1, \boldsymbol{b}_2) $ 化为 $ U $ 的一组标准正交基（ONB）$ C = (\boldsymbol{c}_1, \boldsymbol{c}_2) $，其中：

\begin{equation*}
\boldsymbol{b}_1 := \begin{bmatrix} 1 \\ 1 \\ 1 \end{bmatrix}, \quad \boldsymbol{b}_2 := \begin{bmatrix} -1 \\ 2 \\ 0 \end{bmatrix}
\end{equation*}

\medskip\hrule\medskip

\section*{练习 3.9}
\addcontentsline{toc}{subsection}{练习 3.9}

设 $ n \in \mathbb{N}^* $，并设 $ x_1, \ldots, x_n > 0 $ 为 $ n $ 个正实数，满足 $ x_1 + \cdots + x_n = 1 $。利用柯西-施瓦茨不等式证明：

a.

\begin{equation*}
\sum_{i=1}^n x_i^2 \geqslant \frac{1}{n}
\end{equation*}

b.

\begin{equation*}
\sum_{i=1}^n \frac{1}{x_i} \geqslant n^2
\end{equation*}

提示：考虑 $ \mathbb{R}^n $ 上的点积。然后选择特定的向量 $ \boldsymbol{x}, \boldsymbol{y} \in \mathbb{R}^n $ 并应用柯西-施瓦茨不等式。

\medskip\hrule\medskip

\section*{练习 3.10}
\addcontentsline{toc}{subsection}{练习 3.10}

将向量

\begin{equation*}
\boldsymbol{x}_1 := \begin{bmatrix} 2 \\ 3 \end{bmatrix}, \quad \boldsymbol{x}_2 := \begin{bmatrix} 0 \\ -1 \end{bmatrix}
\end{equation*}

旋转 $ 30^\circ $。
